\section{Die Behauptung}

\textbf{[SETZUNG]} Die Gravitationskonstante $G$ ist in dieser Theorie keine
eigenständige Naturkonstante, sondern eine Einheitenbrücke -- wie $c$,
$\hbar$ und $k_B$ (A080).

Daraus folgt die stärkere und bekanntere Behauptung des Korpus: das Problem der
Quantengravitation stellt sich in dieser Form nicht. Dieses Dokument prüft,
was daran trägt und was nicht.

\section{Warum $G$ eine Brücke ist}

\textbf{[B]} $G$ verbindet eine Masse mit einer Länge. Im natürlichen System
der Theorie, in dem $\tilde T\cdot m = 1$ gilt und Zeit und Masse dieselbe
Größe in zwei Lesarten sind, ist diese Verbindung nicht zusätzlich zu
stiften -- sie liegt bereits in der Dualität.

Das ist dieselbe Aussage wie bei $\hbar$: dort ist die Brücke zwischen Masse
und Zeit gemeint, hier die zwischen Masse und Länge. Beide sind in einem System,
das Zeit, Masse und Länge nicht als unabhängige Einheiten führt, Umrechnungen
und keine Naturgesetze.

\textbf{[SETZUNG]} Operativ bleibt $G$ unverzichtbar. Wer in SI rechnet,
braucht es. Ontologisch ist es sekundär. Das ist genau die Doppelstellung, die
A080 für alle Brückenkonstanten beschreibt, und sie ist keine Abwertung.

\section{Warum sich das QG-Problem in dieser Form nicht stellt}

\textbf{[B]} Das übliche Problem lautet: die Gravitation als Feld quantisieren
und dabei eine renormierbare Theorie erhalten. Es entsteht, weil die Geometrie
als dynamisches Feld behandelt wird, das seinerseits quantisiert werden muss.

In dieser Theorie ist die Geometrie \emph{vorgeordnet} und nicht dynamisch im
selben Sinn: der $T^4/Z_3$ ist gesetzt (A020), die Quantenstruktur entsteht
als sein Spektrum (A030, A070). Es gibt daher nichts, was nachträglich zu
quantisieren wäre.

\textbf{Die ehrliche Bewertung dieser Aussage.} Sie ist kein gelöstes Problem,
sondern ein \emph{aufgelöstes} -- und zwar durch eine Setzung. Wer die
vorgeordnete Geometrie nicht setzt, hat das Problem weiterhin. Die Theorie
beantwortet die Frage nicht, sie stellt sie nicht.

Das ist das \emph{negative} Argument. Ihm steht ein positives zur Seite, das
den Zahlenwert von $G$ aus $\xi$ gewinnt: $G=\xi^2/(4\,m_\text{char})$
(A145). Beide gehören zusammen -- das negative erklärt, warum nicht quantisiert
werden muss, das positive, woher der Wert kommt.

Das ist ein legitimer Zug, aber er darf nicht als Lösung ausgegeben werden. Der
Unterschied ist derselbe wie zwischen einer Herleitung und einer Kalibrierung
(A130): beides zulässig, nur nicht dasselbe.

\section{Was daraus folgt und was nicht}

\textbf{[B]} Es folgt: keine Gravitonen als notwendige Bestandteile, keine
Renormierungsprobleme der Gravitation, keine Singularitäten aus einer
Feldquantisierung.

\textbf{[SETZUNG]} Es folgt \emph{nicht}: dass die klassische
Gravitationsdynamik reproduziert ist. Die A-Serie zeigt an keiner Stelle, dass
sich aus der kompakten Geometrie die beobachtete Gravitationswechselwirkung mit
ihrer Stärke und ihrem Abstandsgesetz ergibt. Das ist eine große offene Stelle
und wird hier nicht kleingeredet.

\section{Prüfskript}

\begin{itemize}
  \item Die Brückenrolle von $G$: dass sich $G$ in natürlichen Einheiten
    genauso absorbieren lässt wie $c$, $\hbar$ und $k_B$, gezeigt an der
    Konsistenz der Einheitengleichungen:
    \texttt{python/A\_Serie\_Skripte/a140\_g\_bruecke.py}
\end{itemize}

\section{Was offen bleibt}

\textbf{Erstens ist die Gravitationsdynamik noch nicht als A-Serien-Dokument
geführt -- sie wird in A142 aufgenommen.} Das vollständige Feldgleichungs-Framework
(Zeitfeld-Lagrange, Skalar-/Fermion-Kopplung, modifizierte Dirac-Gleichung,
Higgs-Kopplung, $\xi$ aus Higgs-Matching) steht im Altkorpus bereit und wird in
A142 eingearbeitet. Was danach noch fehlt: (a) die vollständige Quantisierung
dieser Feldgleichungen, (b) der Übergang zu messbaren Vorhersagen jenseits des
Schwarzschild-Falls, (c) die Verbindung zur Projektionskette (A090).

\textbf{Zweitens ist die Stärke der Wechselwirkung nicht hergeleitet.} $G$ als
Brücke zu führen erklärt nicht, warum die Gravitation um vierzig
Größenordnungen schwächer ist als der Elektromagnetismus.

\textbf{Drittens ist die Auflösung des QG-Problems eine Folge der Setzung.}
Siehe oben; als solche gebucht, nicht als Ergebnis.

\section{Ersetzt}

Dieses Dokument nimmt auf: Dok.~012 (Gravitationskonstante), Dok.~127
(Gravitation), Dok.~163 (Quantengravitation). Eingearbeitet: P40.

\section{Verwendung von KI-Werkzeugen}

Dieses Dokument ist unter Verwendung KI-gestützter Werkzeuge entstanden. Die
Fragestellungen, die Setzungen und alle inhaltlichen Entscheidungen stammen vom
Autor; die Werkzeuge formulieren aus, rechnen nach und erstellen Prüfskripte.
Die Verantwortung für Inhalt und Fehler liegt vollständig beim Autor. Der
ausführliche Wortlaut steht in A010.
